\documentclass{article}
\usepackage{setspace}
\doublespacing
% Language setting
% Replace `english' with e.g. `spanish' to change the document language
\usepackage[english]{babel}

% Set page size and margins
% Replace `letterpaper' with `a4paper' for UK/EU standard size
\usepackage[letterpaper,top=2cm,bottom=2cm,left=3cm,right=3cm,marginparwidth=1.75cm]{geometry}

% Useful packages
\usepackage{amsmath}
\usepackage{graphicx}
\usepackage[colorlinks=true, allcolors=blue]{hyperref}

\title{SNG LaTeX Ken}
\author{Ken Liu}

\begin{document}
\maketitle

\begin{equation}
\theta^{''} = -\frac{3g}{2L}\cos\theta\\
\end{equation}
\begin{equation}
\left\{
\begin{aligned}
x(t) = x_c+t-5\cos\theta\\
y(t)=5\sin\theta
\end{aligned}
\right.
\label{eq:system}
\end{equation}
\noindent The ground is assumed to be parallel to the tangent of the surface of the Earth.

\noindent $\theta \equiv$ the angle that the ladder forms with the ground.

\noindent $\omega \equiv$ the angular velocity of the ladder.

\noindent $\alpha \equiv$ the angular acceleration of the ladder.\\

\noindent and the angle between the ladder and the floor changes with time in accordance with the equation

\begin{equation}
\theta (t) = \arccos \frac{X(t)}{t}
\end{equation}

\noindent where $X(t)$ is given by $X(t) = x_0 + t \cdot V(t)$ and $V(t)$ is given as $1$ m/s.

\noindent Because \[\omega = \lim_{\Delta t -> 0} \frac{\Delta \theta}{\Delta t} = \frac{d\theta}{dt}\] the angular velocity of the falling end of the ladder can be determined by

{\Large
\begin{equation}
\omega(t) = -\frac{V(t) + t\frac{dV}{dt}}{L \sqrt{1 - \frac{(x_0 + tV(t))^2}{L^2}}}
\end{equation}
}
\noindent Furthermore, since 
\[\alpha \equiv \lim_{\Delta t -> 0} \frac{\Delta \omega}{\Delta t} = \frac{d\omega}{dt}\] the angular acceleration of the falling end of the ladder can be given by

{\Large
\begin{equation}
\alpha(t) = - \frac{(x_0+tV(t))(V(t)+t\frac{dV}{dt})^2 + L^2 (1-\frac{(x_0+tV(t))^2}{L^2})(2\frac{dV}{dt}+t\frac{d^2V}{dt^2})}{L^3(1-\frac{(x_0+tV(t))^2}{L^2})^\frac{3}{2}}
\end{equation}
}
\noindent A little more manipulation gives us:

\begin{equation}
v_x(t) = (\frac{k}{m}t+\frac{1}{v_{0x}})^{-1}
\end{equation}

\noindent This gives us the velocity of the ball in the x-direction at time $t$. We also want the position at time $t$, which we derive by integrating:

\begin{equation}
p_x(t) = \int v_x(t)dt = \int(\frac{k}{m}t+\frac{1}{v_{0x}})^{-1} dt = \frac{m}{k} \ln \mid kv_{0x}t+m \mid
\end{equation}

\noindent In the y-direction, equations must take into account both air friction and gravity. We can start the same way as with the x-direction \cite{greenwade93}:
\begin{equation}
F_y = g-kv_y^2 = ma_y = m \frac{dv_y}{dt}
\end{equation}
\bibliographystyle{alpha}
\bibliography{sample}

\end{document}